\section{Overview of Lipid Metabolism}

\begin{learningobjectives}
After completing this chapter, students should be able to:
\begin{enumerate}
    \item Describe the process of fatty acid oxidation ($\beta$-oxidation)
    \item Calculate ATP yield from fatty acid oxidation
    \item Explain fatty acid synthesis and its regulation
    \item Describe ketone body metabolism and its clinical significance
    \item Understand cholesterol synthesis and its regulation
\end{enumerate}
\end{learningobjectives}

\begin{definition}
\textbf{Lipids} serve as:
\begin{itemize}
    \item Major energy storage molecules (triacylglycerols)
    \item Structural components of membranes (phospholipids)
    \item Signaling molecules (steroid hormones, eicosanoids)
    \item Insulation and protection
\end{itemize}
\end{definition}

\section{Fatty Acid Oxidation ($\beta$-Oxidation)}

\subsection{Overview and Location}

\begin{definition}
\textbf{$\beta$-Oxidation} is the process of breaking down fatty acids into acetyl-CoA units for energy production. The name refers to oxidation at the $\beta$-carbon (C3) of the fatty acid.
\end{definition}

\textbf{Location:} Mitochondrial matrix

\textbf{When does it occur?}
\begin{itemize}
    \item Fasting state
    \item Prolonged exercise
    \item Low carbohydrate availability
    \item Diabetes mellitus (uncontrolled)
\end{itemize}

\subsection{Activation of Fatty Acids}

Before oxidation, fatty acids must be activated to acyl-CoA in the cytosol:

\[
\text{Fatty acid} + \text{CoA} + \text{ATP} \xrightarrow{\text{Acyl-CoA synthetase}} \text{Acyl-CoA} + \text{AMP} + \text{PP}_i
\]

\begin{itemize}
    \item Also called fatty acid thiokinase
    \item Located on outer mitochondrial membrane
    \item Costs 2 ATP equivalents (ATP $\rightarrow$ AMP + PP$_i$)
    \item PP$_i$ is hydrolyzed to 2P$_i$ (makes reaction irreversible)
\end{itemize}

\subsection{Carnitine Shuttle}

Long-chain acyl-CoA cannot cross the inner mitochondrial membrane. The \textbf{carnitine shuttle} transports fatty acids into the matrix.

\textbf{Step 1: Formation of Acylcarnitine}
\[
\text{Acyl-CoA} + \text{Carnitine} \xrightarrow{\text{CPT-I}} \text{Acylcarnitine} + \text{CoA}
\]
\begin{itemize}
    \item \textbf{CPT-I} (Carnitine palmitoyltransferase I)
    \item Located on outer mitochondrial membrane
    \item \textbf{Rate-limiting step} of fatty acid oxidation
    \item \textbf{Inhibited by malonyl-CoA} (key regulatory point)
\end{itemize}

\textbf{Step 2: Transport Across Inner Membrane}
\[
\text{Acylcarnitine}_{(\text{cytosol})} \xleftrightarrow{\text{Translocase}} \text{Acylcarnitine}_{(\text{matrix})}
\]
\begin{itemize}
    \item Carnitine-acylcarnitine translocase
    \item Antiport: acylcarnitine in, carnitine out
\end{itemize}

\textbf{Step 3: Regeneration of Acyl-CoA}
\[
\text{Acylcarnitine} + \text{CoA} \xrightarrow{\text{CPT-II}} \text{Acyl-CoA} + \text{Carnitine}
\]
\begin{itemize}
    \item \textbf{CPT-II} on inner membrane (matrix side)
    \item Regenerates acyl-CoA for $\beta$-oxidation
    \item Carnitine returns to cytosol via translocase
\end{itemize}

\begin{tcolorbox}[colback=red!5, colframe=red!70!black, title=\textbf{Clinical Correlation: Carnitine Deficiency}]
\textbf{Primary carnitine deficiency:}
\begin{itemize}
    \item Defect in carnitine transporter
    \item Features: Hypoketotic hypoglycemia, cardiomyopathy, muscle weakness
    \item Treatment: Carnitine supplementation
\end{itemize}

\textbf{CPT-I deficiency:}
\begin{itemize}
    \item Cannot transport fatty acids into mitochondria
    \item Presents with fasting hypoglycemia without ketosis
    \item Hepatomegaly common
\end{itemize}

\textbf{CPT-II deficiency:}
\begin{itemize}
    \item Most common disorder of fatty acid oxidation
    \item Adult form: Exercise-induced muscle pain and rhabdomyolysis
    \item Triggered by prolonged exercise, fasting, cold exposure
\end{itemize}
\end{tcolorbox}

\subsection{The $\beta$-Oxidation Spiral}

Each round of $\beta$-oxidation consists of four reactions that remove a 2-carbon unit as acetyl-CoA.

\textbf{Reaction 1: Oxidation (Acyl-CoA Dehydrogenase)}
\[
\text{Acyl-CoA} + \text{FAD} \rightarrow \text{trans-}\Delta^2\text{-Enoyl-CoA} + \text{FADH}_2
\]
\begin{itemize}
    \item Creates double bond between C2 and C3 (trans configuration)
    \item \textbf{FADH$_2$ produced} $\rightarrow$ 1.5 ATP
    \item Multiple isozymes for different chain lengths:
    \begin{itemize}
        \item VLCAD: Very long chain (C12-C18)
        \item MCAD: Medium chain (C6-C12)
        \item SCAD: Short chain (C4-C6)
    \end{itemize}
\end{itemize}

\textbf{Reaction 2: Hydration (Enoyl-CoA Hydratase)}
\[
\text{trans-}\Delta^2\text{-Enoyl-CoA} + \text{H}_2\text{O} \rightarrow \text{L-3-Hydroxyacyl-CoA}
\]
\begin{itemize}
    \item Adds water across the double bond
    \item Produces L-stereoisomer (important for next step)
\end{itemize}

\textbf{Reaction 3: Oxidation (3-Hydroxyacyl-CoA Dehydrogenase)}
\[
\text{L-3-Hydroxyacyl-CoA} + \text{NAD}^+ \rightarrow \text{3-Ketoacyl-CoA} + \text{NADH} + \text{H}^+
\]
\begin{itemize}
    \item Oxidizes hydroxyl group to ketone
    \item \textbf{NADH produced} $\rightarrow$ 2.5 ATP
\end{itemize}

\textbf{Reaction 4: Thiolysis ($\beta$-Ketothiolase)}
\[
\text{3-Ketoacyl-CoA} + \text{CoA-SH} \rightarrow \text{Acetyl-CoA} + \text{Acyl-CoA}_{(n-2)}
\]
\begin{itemize}
    \item Cleaves between C2 and C3 ($\alpha$ and $\beta$ carbons)
    \item Releases acetyl-CoA
    \item Shortened acyl-CoA re-enters the spiral
\end{itemize}

\begin{tcolorbox}[colback=blue!5, colframe=blue!70!black, title=\textbf{Memory Aid: $\beta$-Oxidation Steps}]
\textbf{``Oh Holy Day, That's Christmas!''}
\begin{itemize}
    \item \textbf{O}xidation (FAD $\rightarrow$ FADH$_2$)
    \item \textbf{H}ydration (add H$_2$O)
    \item \textbf{D}ehydrogenation (NAD$^+$ $\rightarrow$ NADH)
    \item \textbf{T}hiolysis (cleavage with CoA)
\end{itemize}
\end{tcolorbox}

\subsection{Energy Yield from Fatty Acid Oxidation}

\begin{example}
\textbf{Calculate the ATP yield from complete oxidation of palmitate (16:0).}

\textbf{Step 1: Determine number of $\beta$-oxidation cycles}
\begin{itemize}
    \item Palmitate has 16 carbons
    \item Number of cycles = (n/2) - 1 = (16/2) - 1 = \textbf{7 cycles}
\end{itemize}

\textbf{Step 2: Products of $\beta$-oxidation}
\begin{itemize}
    \item Acetyl-CoA: 16/2 = \textbf{8 acetyl-CoA}
    \item FADH$_2$: 7 (one per cycle)
    \item NADH: 7 (one per cycle)
\end{itemize}

\textbf{Step 3: Calculate ATP from each source}

\begin{center}
\begin{tabular}{lcc}
\toprule
\textbf{Source} & \textbf{Amount} & \textbf{ATP} \\
\midrule
Acetyl-CoA (TCA cycle) & 8 $\times$ 10 ATP & 80 ATP \\
FADH$_2$ ($\beta$-oxidation) & 7 $\times$ 1.5 ATP & 10.5 ATP \\
NADH ($\beta$-oxidation) & 7 $\times$ 2.5 ATP & 17.5 ATP \\
\midrule
\textbf{Gross total} & & 108 ATP \\
Activation cost & & $-$2 ATP \\
\midrule
\textbf{Net ATP yield} & & \textbf{106 ATP} \\
\bottomrule
\end{tabular}
\end{center}
\end{example}

\begin{theorem}
\textbf{General Formula for Saturated Fatty Acid Oxidation:}

For a saturated fatty acid with n carbons:
\begin{itemize}
    \item Number of $\beta$-oxidation cycles = (n/2) - 1
    \item Acetyl-CoA produced = n/2
    \item FADH$_2$ produced = (n/2) - 1
    \item NADH from $\beta$-oxidation = (n/2) - 1
\end{itemize}

\textbf{Total ATP} = [n/2 $\times$ 10] + [(n/2 - 1) $\times$ 1.5] + [(n/2 - 1) $\times$ 2.5] - 2
\end{theorem}

\subsection{Oxidation of Unsaturated Fatty Acids}

Unsaturated fatty acids require additional enzymes:

\textbf{1. Enoyl-CoA Isomerase}
\begin{itemize}
    \item Converts \textit{cis}-$\Delta^3$ double bonds to \textit{trans}-$\Delta^2$
    \item Required for monounsaturated fatty acids (e.g., oleate)
    \item No ATP or reducing equivalents involved
\end{itemize}

\textbf{2. 2,4-Dienoyl-CoA Reductase}
\begin{itemize}
    \item Required for polyunsaturated fatty acids
    \item Uses NADPH as reducing agent
    \item Reduces conjugated double bond system
\end{itemize}

\begin{tcolorbox}[colback=yellow!10, colframe=yellow!70!black, title=\textbf{ATP Yield from Unsaturated Fatty Acids}]
Each double bond in a fatty acid results in:
\begin{itemize}
    \item Loss of one FADH$_2$ (1.5 ATP less)
    \item The oxidation step that would create that double bond is bypassed
\end{itemize}

\textbf{Example: Oleate (18:1$\Delta^9$)}
\begin{itemize}
    \item Saturated 18-carbon would yield: 120 - 2 = 118 ATP
    \item Oleate yields: 118 - 1.5 = \textbf{116.5 ATP}
\end{itemize}
\end{tcolorbox}

\subsection{Oxidation of Odd-Chain Fatty Acids}

Odd-chain fatty acids (from bacterial sources, some plants) produce propionyl-CoA in the final thiolysis step.

\textbf{Conversion of Propionyl-CoA to Succinyl-CoA:}

\textbf{Step 1: Carboxylation}
\[
\text{Propionyl-CoA} + \text{CO}_2 + \text{ATP} \xrightarrow{\text{Propionyl-CoA carboxylase}} \text{D-Methylmalonyl-CoA}
\]
\begin{itemize}
    \item Requires \textbf{biotin} as cofactor
    \item Consumes 1 ATP
\end{itemize}

\textbf{Step 2: Racemization}
\[
\text{D-Methylmalonyl-CoA} \xleftrightarrow{\text{Racemase}} \text{L-Methylmalonyl-CoA}
\]

\textbf{Step 3: Isomerization}
\[
\text{L-Methylmalonyl-CoA} \xrightarrow{\text{Methylmalonyl-CoA mutase}} \text{Succinyl-CoA}
\]
\begin{itemize}
    \item Requires \textbf{vitamin B$_{12}$} (adenosylcobalamin)
    \item Succinyl-CoA enters TCA cycle
    \item This is an \textbf{anaplerotic reaction}
\end{itemize}

\begin{tcolorbox}[colback=red!5, colframe=red!70!black, title=\textbf{Clinical Correlation: Methylmalonic Acidemia}]
\textbf{Causes:}
\begin{itemize}
    \item Methylmalonyl-CoA mutase deficiency
    \item Vitamin B$_{12}$ deficiency
\end{itemize}

\textbf{Features:}
\begin{itemize}
    \item Elevated methylmalonic acid in blood and urine
    \item Metabolic acidosis
    \item Developmental delay, lethargy
    \item Distinguishes B$_{12}$ deficiency from folate deficiency
\end{itemize}
\end{tcolorbox}

\section{Fatty Acid Synthesis}

\subsection{Overview}

\begin{definition}
\textbf{Fatty acid synthesis} (lipogenesis) is the process of building fatty acids from acetyl-CoA units. It is essentially the reverse of $\beta$-oxidation but uses different enzymes and occurs in a different cellular location.
\end{definition}

\begin{center}
\begin{tabular}{lcc}
\toprule
\textbf{Feature} & \textbf{$\beta$-Oxidation} & \textbf{Fatty Acid Synthesis} \\
\midrule
Location & Mitochondrial matrix & Cytosol \\
Acyl carrier & CoA & ACP (acyl carrier protein) \\
Reducing agent & FAD, NAD$^+$ & NADPH \\
2-Carbon donor & — & Malonyl-CoA \\
Key enzyme & Multiple enzymes & Fatty acid synthase (FAS) \\
Direction & C
$\rightarrow$ N-terminus & N-terminus $\rightarrow$ C-terminus \\
Energy yield & Produces ATP & Consumes ATP \\
\bottomrule
\end{tabular}
\end{center}

\textbf{Primary sites of fatty acid synthesis:}
\begin{itemize}
    \item \textbf{Liver}: Major site; exports fat as VLDL
    \item \textbf{Adipose tissue}: Storage of synthesized fat
    \item \textbf{Lactating mammary gland}: Milk fat production
    \item \textbf{Brain}: Limited synthesis for myelin
\end{itemize}

\subsection{Transport of Acetyl-CoA to Cytosol}

Acetyl-CoA is produced in mitochondria but fatty acid synthesis occurs in the cytosol. The \textbf{citrate shuttle} transports acetyl-CoA equivalents across the inner mitochondrial membrane.

\textbf{Steps of the Citrate Shuttle:}

\textbf{Step 1: Citrate Formation (Mitochondria)}
\[
\text{Acetyl-CoA} + \text{OAA} \xrightarrow{\text{Citrate synthase}} \text{Citrate}
\]

\textbf{Step 2: Transport}
\begin{itemize}
    \item Citrate exits mitochondria via tricarboxylate transporter
    \item Exchange for malate
\end{itemize}

\textbf{Step 3: Citrate Cleavage (Cytosol)}
\[
\text{Citrate} + \text{ATP} + \text{CoA} \xrightarrow{\text{ATP-citrate lyase}} \text{Acetyl-CoA} + \text{OAA} + \text{ADP} + \text{P}_i
\]
\begin{itemize}
    \item Key regulatory enzyme
    \item Activated by insulin
    \item Inhibited by glucagon
\end{itemize}

\textbf{Step 4: OAA Return to Mitochondria}

Two possible routes:

\textbf{Route A (generates cytosolic NADPH):}
\[
\text{OAA} \xrightarrow{\text{Malate DH}} \text{Malate} \xrightarrow{\text{Malic enzyme}} \text{Pyruvate} + \text{CO}_2 + \text{NADPH}
\]

\textbf{Route B:}
\[
\text{OAA} \xrightarrow{\text{Malate DH}} \text{Malate} \rightarrow \text{(enters mitochondria)}
\]

\begin{tcolorbox}[colback=blue!5, colframe=blue!70!black, title=\textbf{NADPH Sources for Fatty Acid Synthesis}]
Each palmitate synthesis requires \textbf{14 NADPH}:
\begin{enumerate}
    \item \textbf{Pentose phosphate pathway}: Primary source (oxidative phase)
    \item \textbf{Malic enzyme}: Cytosolic; from citrate shuttle
    \item \textbf{Cytosolic isocitrate dehydrogenase}: Minor contribution
\end{enumerate}
\end{tcolorbox}

\subsection{Synthesis of Malonyl-CoA}

\begin{definition}
The committed step of fatty acid synthesis is the carboxylation of acetyl-CoA to form \textbf{malonyl-CoA}, catalyzed by \textbf{acetyl-CoA carboxylase (ACC)}.
\end{definition}

\[
\text{Acetyl-CoA} + \text{CO}_2 + \text{ATP} \xrightarrow{\text{ACC}} \text{Malonyl-CoA} + \text{ADP} + \text{P}_i
\]

\textbf{Properties of ACC:}
\begin{itemize}
    \item Requires \textbf{biotin} as cofactor
    \item Two-step reaction:
    \begin{enumerate}
        \item Biotin carboxylation: Biotin + CO$_2$ + ATP $\rightarrow$ Carboxybiotin
        \item Carboxyl transfer: Carboxybiotin + Acetyl-CoA $\rightarrow$ Malonyl-CoA
    \end{enumerate}
    \item \textbf{Rate-limiting enzyme} of fatty acid synthesis
\end{itemize}

\textbf{Regulation of ACC:}

\begin{center}
\begin{tabular}{ll}
\toprule
\textbf{Activators} & \textbf{Inhibitors} \\
\midrule
Citrate (allosteric) & Palmitoyl-CoA (product inhibition) \\
Insulin (dephosphorylation) & Glucagon (phosphorylation via PKA) \\
High carbohydrate diet & Epinephrine (phosphorylation) \\
 & AMP (via AMPK phosphorylation) \\
\bottomrule
\end{tabular}
\end{center}

\begin{itemize}
    \item \textbf{Active form}: Dephosphorylated polymer
    \item \textbf{Inactive form}: Phosphorylated monomer
    \item Citrate promotes polymerization (activation)
    \item Palmitoyl-CoA promotes depolymerization (inhibition)
\end{itemize}

\subsection{Fatty Acid Synthase (FAS)}

\begin{definition}
\textbf{Fatty acid synthase} is a large, multifunctional enzyme complex that catalyzes all remaining steps of palmitate synthesis. In mammals, it is a homodimer with each monomer containing all seven enzymatic activities.
\end{definition}

\textbf{Domain Organization of FAS:}

\begin{enumerate}
    \item \textbf{Acetyl-CoA-ACP transacylase (AT)}
    \item \textbf{Malonyl-CoA-ACP transacylase (MT)}
    \item \textbf{$\beta$-Ketoacyl-ACP synthase (KS)} — condensing enzyme
    \item \textbf{$\beta$-Ketoacyl-ACP reductase (KR)}
    \item \textbf{$\beta$-Hydroxyacyl-ACP dehydratase (DH)}
    \item \textbf{Enoyl-ACP reductase (ER)}
    \item \textbf{Thioesterase (TE)}
\end{enumerate}

\textbf{Acyl Carrier Protein (ACP):}
\begin{itemize}
    \item Contains phosphopantetheine prosthetic group (from vitamin B$_5$)
    \item Same prosthetic group as CoA
    \item Holds growing fatty acid chain during synthesis
\end{itemize}

\subsection{Reactions of Fatty Acid Synthesis}

\textbf{Loading Reactions:}

\textbf{Step 1: Acetyl Loading}
\[
\text{Acetyl-CoA} + \text{ACP-SH} \xrightarrow{\text{AT}} \text{Acetyl-ACP} + \text{CoA-SH}
\]
Then transferred to KS domain:
\[
\text{Acetyl-ACP} + \text{KS-SH} \rightarrow \text{Acetyl-KS} + \text{ACP-SH}
\]

\textbf{Step 2: Malonyl Loading}
\[
\text{Malonyl-CoA} + \text{ACP-SH} \xrightarrow{\text{MT}} \text{Malonyl-ACP} + \text{CoA-SH}
\]

\textbf{Elongation Cycle (repeated 7 times):}

\textbf{Reaction 1: Condensation ($\beta$-Ketoacyl-ACP Synthase)}
\[
\text{Acetyl-KS} + \text{Malonyl-ACP} \rightarrow \text{Acetoacetyl-ACP} + \text{CO}_2 + \text{KS-SH}
\]
\begin{itemize}
    \item Releases CO$_2$ (same CO$_2$ added by ACC)
    \item Decarboxylation drives reaction forward
    \item Forms $\beta$-ketoacyl intermediate
\end{itemize}

\textbf{Reaction 2: Reduction ($\beta$-Ketoacyl-ACP Reductase)}
\[
\text{Acetoacetyl-ACP} + \text{NADPH} + \text{H}^+ \rightarrow \text{D-$\beta$-Hydroxybutyryl-ACP} + \text{NADP}^+
\]
\begin{itemize}
    \item Reduces ketone to hydroxyl
    \item Uses NADPH (not NADH)
    \item Produces D-isomer (opposite of $\beta$-oxidation)
\end{itemize}

\textbf{Reaction 3: Dehydration ($\beta$-Hydroxyacyl-ACP Dehydratase)}
\[
\text{D-$\beta$-Hydroxybutyryl-ACP} \rightarrow \text{Crotonyl-ACP} + \text{H}_2\text{O}
\]
\begin{itemize}
    \item Removes water
    \item Forms trans double bond (enoyl intermediate)
\end{itemize}

\textbf{Reaction 4: Reduction (Enoyl-ACP Reductase)}
\[
\text{Crotonyl-ACP} + \text{NADPH} + \text{H}^+ \rightarrow \text{Butyryl-ACP} + \text{NADP}^+
\]
\begin{itemize}
    \item Reduces double bond
    \item Second NADPH used per cycle
    \item Product is saturated acyl-ACP (4 carbons)
\end{itemize}

\textbf{Cycle Continuation:}
\begin{itemize}
    \item Butyryl group transferred to KS domain
    \item New malonyl-CoA loaded onto ACP
    \item Cycle repeats, adding 2 carbons each time
    \item After 7 cycles: 16-carbon palmitoyl-ACP
\end{itemize}

\textbf{Termination: Thioesterase}
\[
\text{Palmitoyl-ACP} + \text{H}_2\text{O} \xrightarrow{\text{TE}} \text{Palmitate} + \text{ACP-SH}
\]
\begin{itemize}
    \item Hydrolyzes thioester bond
    \item Releases free palmitate (16:0)
    \item Thioesterase is specific for 16-carbon chains
\end{itemize}

\subsection{Stoichiometry of Palmitate Synthesis}

\begin{tcolorbox}[colback=green!5, colframe=green!70!black, title=\textbf{Overall Reaction for Palmitate Synthesis}]
\[
\text{8 Acetyl-CoA} + \text{7 ATP} + \text{14 NADPH} + \text{14 H}^+ \rightarrow
\]
\[
\text{Palmitate} + \text{8 CoA} + \text{7 ADP} + \text{7 P}_i + \text{14 NADP}^+ + \text{6 H}_2\text{O}
\]

\textbf{Accounting:}
\begin{itemize}
    \item 8 acetyl-CoA provides 16 carbons
    \item 7 ATP used to make 7 malonyl-CoA
    \item 14 NADPH: 2 per elongation cycle $\times$ 7 cycles
    \item 7 CO$_2$ added then removed (net zero)
\end{itemize}
\end{tcolorbox}

\subsection{Elongation and Desaturation}

\textbf{Fatty Acid Elongation:}

Palmitate can be elongated to longer fatty acids:

\begin{itemize}
    \item \textbf{Endoplasmic reticulum}: Primary site
    \begin{itemize}
        \item Uses malonyl-CoA as 2-carbon donor
        \item Uses NADPH as reductant
        \item Similar reactions to FAS but separate enzymes
    \end{itemize}
    \item \textbf{Mitochondria}: Minor pathway
    \begin{itemize}
        \item Uses acetyl-CoA directly
        \item Reverse of $\beta$-oxidation
    \end{itemize}
\end{itemize}

\textbf{Fatty Acid Desaturation:}

\begin{definition}
\textbf{Desaturases} introduce double bonds into fatty acid chains. Mammals have $\Delta$9, $\Delta$6, $\Delta$5, and $\Delta$4 desaturases but \textbf{cannot} introduce double bonds beyond C-9.
\end{definition}

\textbf{$\Delta$9-Desaturase (Stearoyl-CoA Desaturase):}
\[
\text{Stearoyl-CoA} + \text{NADH} + \text{H}^+ + \text{O}_2 \rightarrow \text{Oleoyl-CoA} + \text{NAD}^+ + 2\text{H}_2\text{O}
\]

\begin{itemize}
    \item Located in ER membrane
    \item Requires cytochrome b$_5$ and cytochrome b$_5$ reductase
    \item Mixed-function oxidase (uses O$_2$)
    \item Creates cis double bond at $\Delta$9 position
    \item Converts stearate (18:0) to oleate (18:1$\Delta$9)
\end{itemize}

\begin{tcolorbox}[colback=red!5, colframe=red!70!black, title=\textbf{Essential Fatty Acids}]
Because humans lack $\Delta$12 and $\Delta$15 desaturases, we cannot synthesize:

\textbf{$\omega$-6 family:}
\begin{itemize}
    \item Linoleic acid (18:2$\Delta$9,12) — essential
    \item Arachidonic acid (20:4$\Delta$5,8,11,14) — conditionally essential
\end{itemize}

\textbf{$\omega$-3 family:}
\begin{itemize}
    \item $\alpha$-Linolenic acid (18:3$\Delta$9,12,15) — essential
    \item EPA (20:5) and DHA (22:6) — conditionally essential
\end{itemize}

These must be obtained from the diet!
\end{tcolorbox}

\subsection{Regulation of Fatty Acid Synthesis}

\textbf{Short-term Regulation (minutes):}

\begin{enumerate}
    \item \textbf{Allosteric regulation of ACC:}
    \begin{itemize}
        \item Citrate activates (indicates abundant acetyl-CoA)
        \item Palmitoyl-CoA inhibits (product feedback)
    \end{itemize}
    
    \item \textbf{Covalent modification of ACC:}
    \begin{itemize}
        \item \textbf{Insulin}: Activates phosphatase $\rightarrow$ dephosphorylates ACC $\rightarrow$ active
        \item \textbf{Glucagon/Epinephrine}: Activates PKA $\rightarrow$ phosphorylates ACC $\rightarrow$ inactive
        \item \textbf{AMPK}: Phosphorylates ACC when ATP is low $\rightarrow$ inactive
    \end{itemize}
\end{enumerate}

\textbf{Long-term Regulation (hours to days):}

\begin{enumerate}
    \item \textbf{Transcriptional control:}
    \begin{itemize}
        \item \textbf{SREBP-1c}: Activated by insulin; induces ACC, FAS, ATP-citrate lyase
        \item \textbf{ChREBP}: Activated by glucose; induces lipogenic enzymes
    \end{itemize}
    
    \item \textbf{Nutritional state:}
    \begin{itemize}
        \item High carbohydrate diet: Increases lipogenic enzyme synthesis
        \item Fasting/starvation: Decreases lipogenic enzyme synthesis
        \item High fat diet: Decreases lipogenic enzymes
    \end{itemize}
\end{enumerate}

\section{Ketone Body Metabolism}

\subsection{Overview of Ketogenesis}

\begin{definition}
\textbf{Ketone bodies} are water-soluble molecules produced in the liver from acetyl-CoA during periods of low carbohydrate availability. They serve as an alternative fuel source for extrahepatic tissues, especially the brain.
\end{definition}

\textbf{The three ketone bodies:}
\begin{enumerate}
    \item \textbf{Acetoacetate} — primary ketone body
    \item \textbf{$\beta$-Hydroxybutyrate} (D-3-hydroxybutyrate) — most abundant in blood
    \item \textbf{Acetone} — produced by spontaneous decarboxylation; exhaled
\end{enumerate}

\textbf{Conditions favoring ketogenesis:}
\begin{itemize}
    \item Prolonged fasting/starvation
    \item Uncontrolled diabetes mellitus
    \item Very low carbohydrate (ketogenic) diets
    \item Alcoholism
    \item High-intensity prolonged exercise
\end{itemize}

\begin{tcolorbox}[colback=blue!5, colframe=blue!70!black, title=\textbf{Why Ketone Bodies?}]
\textbf{Problem:} During fasting, the brain needs fuel but:
\begin{itemize}
    \item Cannot use fatty acids (don't cross blood-brain barrier efficiently)
    \item Glucose supply is limited (gluconeogenesis is slow)
\end{itemize}

\textbf{Solution:} Liver converts fatty acids to ketone bodies, which:
\begin{itemize}
    \item Are water-soluble (cross blood-brain barrier)
    \item Spare glucose for tissues that absolutely require it (RBCs)
    \item Reduce protein breakdown for gluconeogenesis
\end{itemize}
\end{tcolorbox}

\subsection{Pathway of Ketogenesis}

\textbf{Location:} Mitochondrial matrix of \textbf{hepatocytes only}

\textbf{Substrate:} Acetyl-CoA from $\beta$-oxidation of fatty acids

\textbf{Step 1: Formation of Acetoacetyl-CoA}
\[
2 \text{ Acetyl-CoA} \xrightarrow{\text{Thiolase}} \text{Acetoacetyl-CoA} + \text{CoA-SH}
\]
\begin{itemize}
    \item Reverse of the last step of $\beta$-oxidation
    \item Condensation of two acetyl-CoA molecules
\end{itemize}

\textbf{Step 2: Formation of HMG-CoA}
\[
\text{Acetoacetyl-CoA} + \text{Acetyl-CoA} + \text{H}_2\text{O} \xrightarrow{\text{HMG-CoA synthase}} \text{HMG-CoA} + \text{CoA-SH}
\]
\begin{itemize}
    \item \textbf{HMG-CoA synthase}: Rate-limiting enzyme of ketogenesis
    \item HMG-CoA = 3-hydroxy-3-methylglutaryl-CoA
    \item This is the \textbf{mitochondrial} HMG-CoA synthase (distinct from cytosolic enzyme in cholesterol synthesis)
\end{itemize}

\textbf{Step 3: Cleavage to Acetoacetate}
\[
\text{HMG-CoA} \xrightarrow{\text{HMG-CoA lyase}} \text{Acetoacetate} + \text{Acetyl-CoA}
\]
\begin{itemize}
    \item Releases the first ketone body (acetoacetate)
    \item Acetyl-CoA is recycled
\end{itemize}

\textbf{Step 4a: Reduction to $\beta$-Hydroxybutyrate}
\[
\text{Acetoacetate} + \text{NADH} + \text{H}^+ \xleftrightarrow{\beta\text{-hydroxybutyrate DH}} \beta\text{-Hydroxybutyrate} + \text{NAD}^+
\]
\begin{itemize}
    \item Reversible reaction
    \item Ratio of $\beta$-hydroxybutyrate to acetoacetate reflects mitochondrial NADH/NAD$^+$ ratio
    \item Typically 3:1 under normal conditions; higher in severe ketosis
\end{itemize}

\textbf{Step 4b: Spontaneous Decarboxylation to Acetone}
\[
\text{Acetoacetate} \xrightarrow{\text{spontaneous}} \text{Acetone} + \text{CO}_2
\]
\begin{itemize}
    \item Non-enzymatic reaction
    \item Acetone is volatile; exhaled through lungs
    \item Causes "fruity breath" odor in diabetic ketoacidosis
\end{itemize}

\begin{figure}[h]
\centering
\begin{tcolorbox}[colback=white, colframe=black, width=0.9\textwidth]
\textbf{Summary of Ketogenesis Pathway:}
\[
\text{Fatty Acids} \xrightarrow{\beta\text{-oxidation}} \text{Acetyl-CoA} \xrightarrow{\text{Thiolase}} \text{Acetoacetyl-CoA}
\]
\[
\downarrow \text{ HMG-CoA synthase}
\]
\[
\text{HMG-CoA} \xrightarrow{\text{HMG-CoA lyase}} \text{Acetoacetate}
\]
\[
\swarrow \qquad \qquad \searrow
\]
\[
\beta\text{-Hydroxybutyrate} \qquad \text{Acetone}
\]
\end{tcolorbox}
\end{figure}

\subsection{Ketolysis (Ketone Body Utilization)}

\begin{definition}
\textbf{Ketolysis} is the process by which extrahepatic tissues convert ketone bodies back to acetyl-CoA for energy production via the TCA cycle.
\end{definition}

\textbf{Tissues that use ketone bodies:}
\begin{itemize}
    \item \textbf{Brain}: Major consumer during prolonged fasting (up to 75\% of energy needs)
    \item \textbf{Heart}: Prefers ketone bodies over glucose
    \item \textbf{Skeletal muscle}: Uses ketones during fasting
    \item \textbf{Kidney cortex}: Active ketone consumer
\end{itemize}

\textbf{Tissues that CANNOT use ketone bodies:}
\begin{itemize}
    \item \textbf{Liver}: Lacks succinyl-CoA:acetoacetate CoA transferase (SCOT)
    \item \textbf{Red blood cells}: No mitochondria
\end{itemize}

\textbf{Steps of Ketolysis:}

\textbf{Step 1: Oxidation of $\beta$-Hydroxybutyrate}
\[
\beta\text{-Hydroxybutyrate} + \text{NAD}^+ \xrightarrow{\beta\text{-hydroxybutyrate DH}} \text{Acetoacetate} + \text{NADH}
\]

\textbf{Step 2: Activation of Acetoacetate}
\[
\text{Acetoacetate} + \text{Succinyl-CoA} \xrightarrow{\text{SCOT}} \text{Acetoacetyl-CoA} + \text{Succinate}
\]
\begin{itemize}
    \item \textbf{SCOT} = Succinyl-CoA:acetoacetate CoA transferase (thiophorase)
    \item Key enzyme absent in liver (prevents futile cycle)
    \item Uses succinyl-CoA from TCA cycle as CoA donor
\end{itemize}

\textbf{Step 3: Thiolytic Cleavage}
\[
\text{Acetoacetyl-CoA} + \text{CoA-SH} \xrightarrow{\text{Thiolase}} 2 \text{ Acetyl-CoA}
\]

\textbf{Step 4: Entry into TCA Cycle}
\[
2 \text{ Acetyl-CoA} \xrightarrow{\text{TCA cycle}} \text{ATP production}
\]

\begin{tcolorbox}[colback=green!5, colframe=green!70!black, title=\textbf{ATP Yield from Ketone Bodies}]
\textbf{From $\beta$-Hydroxybutyrate:}
\begin{itemize}
    \item 1 NADH (from $\beta$-HB dehydrogenase) = 2.5 ATP
    \item 2 Acetyl-CoA $\times$ 10 ATP = 20 ATP
    \item Minus 1 GTP equivalent (succinyl-CoA used for activation) = -1 ATP
    \item \textbf{Total: 21.5 ATP}
\end{itemize}

\textbf{From Acetoacetate:}
\begin{itemize}
    \item 2 Acetyl-CoA $\times$ 10 ATP = 20 ATP
    \item Minus 1 GTP equivalent = -1 ATP
    \item \textbf{Total: 19 ATP}
\end{itemize}
\end{tcolorbox}

\subsection{Regulation of Ketogenesis}

\textbf{Factors promoting ketogenesis:}

\begin{enumerate}
    \item \textbf{Increased fatty acid supply to liver:}
    \begin{itemize}
        \item Low insulin $\rightarrow$ increased lipolysis in adipose tissue
        \item High glucagon/epinephrine $\rightarrow$ activated hormone-sensitive lipase
    \end{itemize}
    
    \item \textbf{Increased $\beta$-oxidation:}
    \begin{itemize}
        \item Low malonyl-CoA (glucagon effect) $\rightarrow$ active CPT-I
        \item More fatty acids enter mitochondria
    \end{itemize}
    
    \item \textbf{Decreased TCA cycle activity:}
    \begin{itemize}
        \item Oxaloacetate diverted to gluconeogenesis
        \item Acetyl-CoA accumulates and enters ketogenesis
    \end{itemize}
    
    \item \textbf{High NADH/NAD$^+$ ratio:}
    \begin{itemize}
        \item From excessive $\beta$-oxidation
        \item Inhibits isocitrate dehydrogenase and $\alpha$-ketoglutarate dehydrogenase
        \item Slows TCA cycle, promoting ketone body formation
    \end{itemize}
\end{enumerate}

\begin{tcolorbox}[colback=red!5, colframe=red!70!black, title=\textbf{Clinical Correlation: Diabetic Ketoacidosis (DKA)}]
\textbf{Pathophysiology:}
\begin{itemize}
    \item Insulin deficiency (Type 1 DM) or severe insulin resistance
    \item Unopposed glucagon action
    \item Massive lipolysis $\rightarrow$ excess fatty acids to liver
    \item Overwhelming ketone body production
\end{itemize}

\textbf{Clinical Features:}
\begin{itemize}
    \item Hyperglycemia ($>$250 mg/dL)
    \item Metabolic acidosis (pH $<$ 7.3, bicarbonate $<$ 18 mEq/L)
    \item Ketonemia and ketonuria
    \item Fruity breath (acetone)
    \item Kussmaul breathing (deep, rapid respirations)
    \item Dehydration, altered mental status
\end{itemize}

\textbf{Treatment:}
\begin{itemize}
    \item IV fluids (correct dehydration)
    \item IV insulin (stop ketogenesis)
    \item Potassium replacement (insulin drives K$^+$ into cells)
    \item Monitor and correct electrolytes
\end{itemize}
\end{tcolorbox}

\subsection{Starvation Ketosis vs. Diabetic Ketoacidosis}

\begin{center}
\begin{tabular}{lcc}
\toprule
\textbf{Feature} & \textbf{Starvation Ketosis} & \textbf{DKA} \\
\midrule
Insulin level & Low but present & Absent or very low \\
Glucagon level & Elevated & Very elevated \\
Blood glucose & Low to normal & High ($>$250 mg/dL) \\
Ketone levels & Moderate (2-5 mM) & Very high ($>$10 mM) \\
pH & Mildly decreased & Severely decreased ($<$7.3) \\
Severity & Self-limited & Life-threatening \\
\bottomrule
\end{tabular}
\end{center}

